\subsection{Cyclic Groups}

\begin{definition}
    Given a group $G$ and an element $g \in G$, the cyclic
    subgroup generated by $g$ is \(\langle g \rangle= \left\{ g^m \mid m \in\Z \right\}
    = \left\{ \ldots, g^{-2}, g^{-1}, e, g, g^2, \ldots \right\}\)\\
    We say $G$ is a cyclic group if there is a $g\in G$ s.t. $G=\langle g \rangle$\\
    ord$(g)=\mid\langle g \rangle \mid$
\end{definition}

\subsubsection*{Cyclic subgroups of $\Z$}
\begin{example}
    \phantom{text}

    \begin{itemize}
        \item $\ang{2}= \left\{ 2m \mid m \in\Z \right\}
        = \left\{ \ldots, -4, -2, 0, 2, 4, \ldots \right\}
        = 2\Z$
        \item $\ang{n}= \left\{ nm \mid m \in\Z \right\}
        = \left\{ \ldots, -2n, -n, 0, n, 2n, \ldots \right\}
        = n\Z$
    \end{itemize}
\end{example}
\subsubsection*{Cyclic subgroups of $\Q^{\times}$}
\begin{example}
    \phantom{text}

    \begin{itemize}
        \item $\ang{-1}=\left\{ -1,1 \right\}$
        \item $\ang{2}=\left\{ 2^m \mid m\in\Z \right\}$
    \end{itemize}
\end{example}

\subsubsection*{Cyclic subgroups of $\C^{\times}$}
\begin{example}
    \phantom{text}

    $\ang{i}=\left\{ i,-1,-i,1 \right\}$
\end{example}

\subsubsection*{Cyclic subgroups of $\Z_{12}$}
\begin{example}
    \phantom{text}

    \begin{itemize}
        \item $\Z_{12}=\ang{1}=\left\{ 0,1,2,3,4,5,6,7,8,9,10,11 \right\}$
        \item $\ang{2}=\left\{ 0,2,4,6,8,10 \right\}$
        \item $\ang{4}=\left\{ 0,4,8 \right\}$
        \item $\ang{3}=\left\{ 0,3,6,9 \right\}$
    \end{itemize}
\end{example}

\subsubsection*{Cyclic subgroups of $U_9=\left\{ 1,2,4,5,7,8 \right\}$}
\begin{example}
    \phantom{text}

    \begin{itemize}
        \item $\ang{1}=\left\{ 1 \right\}$
        \item $\ang{8}=\left\{ 8,64 \right\}=\left\{ 1,8 \right\}$
        \item $\ang{4}=\left\{ 4,16,64 \right\}=\left\{ 1,4,7 \right\}$
        \item $\ang{2}=\left\{ 2,4,8,16,32,64 \right\}=\left\{ 1,2,4,5,7,8 \right\}=U_9$
    \end{itemize}
\end{example}

\subsubsection*{Cyclic subgroups of $D_n$}
\begin{example}
    \phantom{text}

    \begin{itemize}
        \item $\ang{r}=\left\{ e,r,r^2,\ldots,r^{n-1} \right\}=Z_n\footnote{not $\Z$}=C_n\footnote{Multiplicative cyclic group of order n}$
        \item $\ang{s}=\left\{ e,s \right\}$
        \item $\ang{sr}=\left\{ e,sr \right\}$
    \end{itemize}
\end{example}
\subsubsection*{Cyclic subgroups of $S_5$}
\begin{example}
    \phantom{text}

    $\ang{(12)(345)}\footnotemark
    =\left\{ (1),(12)(345),(354),(12),(345),(12)(354)\right\}$
    \footnotetext{By \Cref{thm:orderlcm} the lcm(2,3)=6=order and thus this cyclic group will have 6 elements}
\end{example}

\begin{theorem}
    \label[theorem]{thm:cyclicgcd}
    Suppose $g \in G$ with ord$(g)=n$ and let $m\in\N$ then, $\ang{g^m}=\ang{g^{\t{gcd}(m,n)}}$
\end{theorem}
\begin{proof}
    Set $d=\t{gcd}(m,n)$.

    Take $x \in \ang{g^m}\implies x=(g^m)^k=g^{mk}$ for some $k \in\Z$

    Also, $m=dq$ for some $q \in\Z$ since $d$ is a divisor of $m$

    $x=g^{dqk}=(g^d)^{qk}\in \ang{g^d}$ so $\ang{g} \subseteq \ang{g^d}$

    By \Cref{lem:bezout} there exits $r,s \in\Z$ s.t. $mr+ns=d$

    Take $x \in \ang{g^d} \implies x=g^{dk}$ for some $k \in\Z$

    $g^{dk}=g^{k(mr+ns)}=g^{mkr}(g^{n})^{ks}=g^{mkr}\in\ang{g^m}$ so $\ang{g^d}\subseteq\ang{g^m}$

    Thus $\ang{g^m}=\ang{g^d}$
\end{proof}

\begin{theorem}[Fundamental Theorem of Cyclic Groups]
    \label[theorem]{thm:fundamentalthmofcyclic}
    \phantom{text}

    \begin{enumerate}[I.]
        \item Every subgroup of a cyclic group is cyclic
        \item If $G$ is a cyclic group of order $n$ and $d\mid n$ then
        there is a unique subgroup $H \leq G$ with $|H|=d$
    \end{enumerate}
\end{theorem}

\begin{enumerate}[I.]
    \item \begin{proof}
        Suppose $G=\ang{g}$ and $H \leq G$

        Case 1: $H=\left\{ e \right\}=\ang{e}$ \checkmark

        Case 2: $H \neq \{e\}$.
        Thus, $H=\{\ldots, g^{-m_2},g^{-m_1},e,g^{m_1},g^{m_2},\ldots\}$ s.t. $m_i\in\N$

        By the Well Ordering Principle, there exists a minimun $m_i$ call this $m$

        Since $g^m\in H \implies (g^m)^n \in H\;\;\forall n\in\Z \implies \ang{g^m}\subseteq H$

        Take $h\in H$, write $h=g^k$ for $k \in \Z$

        By the division algorithm (Thm. \ref{thm:divalgo}), there exits a $q,r \in \Z$ s.t.\\
        $k=mq+r$ with $0\leq r<m$

        $g^r=(g^{k})(g^{-mq})$ thus since $g^k\in H$ and $g^{-mq}\in \ang{g^m}\subseteq H \implies g^r\in H$

        Since $0\leq r < m$ and $m$ is the minimun $\implies r=0$

        Thus $h=(g^m)^q\in\ang{g^m} \implies H \subseteq \ang{g^m}$

        $\implies H=\ang{g^m}\implies H$ is cyclic
    \end{proof}
    \item \begin{proof}
        Assume $G$ is a cyclic group of order $n$ and $d\mid n$

        Let $G=\ang{g}$

        Existence: Consider $\ang{g^{\frac{n}{d}}}$ this has $d$ elements so \checkmark

        Uniqueness: Let $H \leq G$, $|H|=d$ by \Cref{thm:fundamentalthmofcyclic}.I 
        
        $\implies H=\ang{g^m}$ by \Cref{thm:cyclicgcd} $\implies H=\ang{g^{\t{gcd}(m,n)}}$

        Thus $|H|=d=\frac{n}{\t{gcd}(m,n)} \implies \t{gcd}(m,n)=\frac{n}{d} \implies H=\ang{g^\frac{n}{d}}$ which unique describes a group
    \end{proof}
\end{enumerate}

\begin{definition}[Euler's Totient Function]
    \phantom{text}\\
    For $n\in\N$, define $\varphi(n)=\Big|\left\{ 1\leq m \leq n \mid \t{gcd}(m,n)=1 \right\}\Big|$
\end{definition}

\begin{theorem}
    \label[theorem]{thm:cyclictoitent}
    \phantom{text}\\
    Suppose $G$ is a cyclic group of order $n$ and $d\mid n$
    then the number of elements of $G$ of order $d$ is $\varphi(d)$
\end{theorem}

\begin{proof}
    Assume $G$ is cyclic group of order $n$ and $d\mid n$

    Let $H\leq G$ be the unique subgroup s.t. $|H|=d$
    $\impliedby$ \Cref{thm:fundamentalthmofcyclic}.II

    Let $H=\ang{h} \impliedby$ \Cref{thm:fundamentalthmofcyclic}.I

    Let $x\in G$ have order $d$
    
    Consider: $\ang{x}$ since $x$ is order $d \implies |\ang{x}|=d$

    Since H is unique $\implies \ang{h}=\ang{x} \implies x=h^m$ where gcd$(m,d)=1 \impliedby$ \Cref{thm:cyclicgcd}

    The number of choices of $m$ is $\varphi(d)$
\end{proof}

\newpage
\subsubsection{Exercises}
\begin{example}
    Find all generators of $\Z_{15} \t{ and } Z_{12}$

    By \Cref{thm:cyclicgcd} this implies all the numbers that are relatively prime to 12 and 15

    Thus, the generators for $\Z_{15}$ are 1,2,4,7,8,11,13,14

    Thus, the generators for $\Z_{12}$ are 1,5,7,11
    
\end{example}

\begin{example}
    Show that $\Q^\times_{\geq0}$ is not cyclic

    Assume $\Q^\times_{\geq0}$ is cyclic thus there exists
    a $q=\frac{a}{b}\in\Q^\times_{\geq0}$ s.t. $\Q^\times_{\geq0}=\ang{q}$

    Consider: $p$ a prime bigger than $a,b$

    Let $\frac{1}{p}=\frac{a^m}{b^m}$

    Case 1: $m \geq 1$

    $\implies b^m=pa^m \implies p \mid b^m \contr \t{since we picked p to be bigger than } a,b$

    Case 2: $m \leq -1$
    
    $\implies a^m=pb^m \implies p \mid a^m \contr$
\end{example}